\section{Related Work}
\label{sec:related}

\paragraph{Context-conditioned and meta-reinforcement learning.}
Meta-RL and context-conditioned control address hidden task or dynamics variation by learning policies that adapt from experience. Gradient-based meta-learning and learned dynamics adaptation can tune policies or models after observing a small amount of data~\cite{finn2017model,nagabandi2019learning}, while latent-context methods infer a task embedding from recent transitions and condition the policy on that embedding~\cite{rakelly2019efficient,zintgraf2021varibad}. \PTA{} uses this latter idea in a deliberately ordinary form: a short trace is encoded into a latent context, and a residual policy conditions on that latent. The difference is not a new inference architecture. Instead, we separate the data used to infer context from the data produced by the task rollout, then evaluate a complication often hidden by standard adaptation benchmarks: in contact-rich materials, the experience used for inference is also an intervention that changes the state distribution.

\paragraph{Active system identification and information-gathering control.}
Online adaptation has also been studied as system identification, where policies or models infer low-dimensional dynamics parameters during deployment. Model-based meta-learning and online adaptation methods update dynamics estimates from recent interaction~\cite{nagabandi2018deep,clavera2019modelbased}, and sim-to-real pipelines often combine broad training variation with online estimation of physical parameters or actuator properties~\cite{yu2017preparing,peng2018sim2real}. These approaches usually rely on informative transitions that arise during normal task execution, or on identification signals whose cost is not the main experimental object. \PTA{} makes the information-gathering phase explicit but minimal: the probe is fixed, short, and open-loop, while the learned encoder and downstream controller are standard. This lets us ask a diagnostic question that is easy to obscure in more capable adaptive systems: when the information-gathering action changes the material state before the task begins, does the learned context outweigh the physical cost of acquiring it?

\paragraph{Embodied probing and physical exploration.}
Robotics has long treated perception as an active process in which actions are selected to reveal otherwise hidden physical information. Early active vision~\cite{bajcsy1988active} and haptic exploratory procedures~\cite{lederman1987hand} established that contact motions can be designed around the property to be sensed. Modern tactile manipulation uses exploratory contact to estimate local geometry, contact conditions, regrasp affordances, and object pose~\cite{calandra2018more,xu2013tactile,bauza2020tactile,yu2016more,bauza2018data}. This literature motivates probing as a physical measurement operation, but many studied settings involve rigid objects or local contact properties where the measurement action need not substantially degrade the downstream task state. Our setting is intentionally less benign: the same push that reveals stiffness, cohesion, or plasticity can also leave grooves, compact material, or otherwise alter the target configuration. The recoverability of that disturbance is therefore an interpretive lens for our results, not an independently identified law: it is consistent with the one positive elastoplastic split and the four negative splits, but further material families and probe designs would be needed to validate it as a boundary condition.

\paragraph{Deformable-material benchmarks as hidden-dynamics stress tests.}
Recent simulators and benchmarks make it practical to evaluate control under diverse soft-body and granular dynamics. Genesis supports unified rigid, MPM, SPH, and FEM simulation~\cite{genesis2024}; PlasticineLab and DiffSkill study skill learning in differentiable MPM environments~\cite{huang2022plasticinelab,lin2022diffskill}; RoboCraft and RoboCook examine tool use with deformable materials~\cite{shi2022robocraft,shi2023robocook}; and SoftGym provides deformable manipulation tasks across cloth, rope, and fluid-like domains~\cite{lin2021softgym}. These benchmarks emphasize deformable control, planning, or differentiable simulation, whereas our benchmark fixes a single edge-pushing task and varies the hidden material regime to stress-test context acquisition under controlled train/test splits. That design makes the result interpretable as a diagnostic rather than as a leaderboard claim: a sand-trained policy with explicit probing improves transfer only on the recoverable elastoplastic split, while the same mechanism is worse on four irreversible or partly irreversible splits. Our novelty claim is therefore not architectural; it is the controlled diagnostic evaluation of a minimal active-context-acquisition mechanism across material regimes where the information-gathering action can itself alter the task state.
